\section{Implementation Plan}
The implementation is organized into data preparation, patch extraction, model training, inference and evaluation. Table~\ref{tab:implementation-status} summarizes what has been completed at the interim stage and what remains for the final submission.

\begin{table}[H]
\centering
\scriptsize
\renewcommand{\arraystretch}{1.25}

\caption{Implementation status at the interim submission.}

\begin{tabularx}{\columnwidth}{|
>{\raggedright\arraybackslash}p{1.8cm}|
>{\raggedright\arraybackslash}p{2cm}|
X|
>{\centering\arraybackslash}p{1.4cm}|}

\hline


\textbf{Module} &
\textbf{Component} &
\textbf{Functionality} &
\textbf{Status} \\ \hline

Data Preprocessing &
Dataset Indexing &
Parse MVTec 3D-AD structure, splits and metadata &
\cellcolor{implgreen}Implemented \\ \cline{2-4}

&
Depth Preprocessing &
Extract depth channel, mask invalid values, crop, resize and normalize depth maps &
\cellcolor{implgreen}Implemented \\ \hline

Patch Extraction &
Patch Generation &
Generate overlapping $32 \times 32$ patches &
\cellcolor{implgreen}Implemented \\ \cline{2-4}

&
Score Aggregation &
Aggregate patch scores into image space &
\cellcolor{implgreen}Implemented \\ \hline

Classical Baseline &
Current Baseline &
Use raw depth patches with PCA and One-Class SVM &
\cellcolor{implgreen}Implemented \\ \cline{2-4}

&
Test Inference &
Generate image scores and anomaly maps &
\cellcolor{implgreen}Implemented \\ \cline{2-4}

&
Preliminary RGB Analysis &
Analyze whether RGB contains useful anomaly information &
\cellcolor{implgreen}Implemented \\ \cline{2-4}

&
Future Extensions &
Explore RGB-enhanced or multimodal representations &
\cellcolor{planblue}Planned \\ \hline

Deep Learning Model &
Autoencoder Design &
Design convolutional autoencoder &
\cellcolor{progYellow}In Progress \\ \cline{2-4}

&
Autoencoder Training &
Train reconstruction-based model &
\cellcolor{planblue}Planned \\ \hline

Evaluation \& Visualization &
Classical Evaluation &
Compute AUROC and AP for the classical baseline &
\cellcolor{implgreen}Implemented \\ \cline{2-4}

&
Classical Heatmaps &
Generate anomaly heatmaps for the classical baseline &
\cellcolor{implgreen}Implemented \\ \cline{2-4}

&
Autoencoder Evaluation &
Compute AUROC, AP and heatmaps for the autoencoder &
\cellcolor{planblue}Planned \\ \cline{2-4}

&
Comparative Evaluation &
Compare classical and deep learning results &
\cellcolor{planblue}Planned \\ \hline

\end{tabularx}

\label{tab:implementation-status}
\end{table}

\subsection*{Repository organization and reproducibility}
The implementation is organized as a reproducible repository rather than as isolated scripts. Shared parameters are stored in \texttt{configs/}, reusable code lives in \texttt{src/}, runnable entry points are placed in \texttt{scripts/}, visual evidence is saved in \texttt{fig/}, and trained models, logs and metrics are stored under \texttt{outputs/}. The main modules are separated by responsibility: \texttt{src/data/} handles indexing, preprocessing and patch extraction; \texttt{src/features/} contains feature construction; \texttt{src/models/} contains anomaly models; \texttt{src/training/} and \texttt{src/inference/} implement model execution; and \texttt{src/evaluation/} contains metrics and visualization utilities.

The current pipeline can be reproduced through documented commands:
\begin{lstlisting}[basicstyle=\ttfamily\scriptsize,frame=single]
conda env create -f environment.yml
conda activate mvtec-3d-ad
pip install -e .
python scripts/prepare_data.py
python scripts/train_classical.py
python scripts/infer_classical.py --split test
python -m unittest discover tests
\end{lstlisting}

This structure makes the interim results auditable: the reported metrics come from saved CSV files, the heatmaps are generated by the inference pipeline, and the classical models are stored per category with their scaler, PCA transform and feature metadata. The next implementation priority is to complete the autoencoder training/inference path and run a controlled modality ablation before the final comparison.
